\section{Reduction parameters and finite Batalin--Vilkovisky complexes}

\begin{proposition}[Conformal obstruction to parameter reflection]
\label{prop:cbc26-conformal-reflection}
Let $t\in\mathbb C^*$ and put $c(t)=13-6(t+t^{-1})$.
A conformal vertex algebra isomorphism between the universal Virasoro algebras at charges $c(t)$ and $c(-t)$ exists if and only if $t=\pm i$.
\end{proposition}
\begin{proof}
A conformal isomorphism preserves the vacuum and conformal vector $T$.
It therefore preserves $T_{(3)}T=(c/2)\mathbf1$, so the two charges must agree.
Since their sum is $26$, equality is equivalent to $c(t)=13$, or $t+t^{-1}=0$.
This gives $t=\pm i$.
Conversely both universal algebras then have charge $13$, and the map sending their conformal generators to one another is an isomorphism by their identical universal presentation.
\end{proof}

At $t=-2/3$, reflection gives $c(t)=26$ and $c(-t)=0$, whereas the reciprocal partner $t^{-1}=-3/2$ has charge $26$.
This distinguishes reflection from conformal Feigin--Frenkel duality by an actual vertex product.
It does not exclude a curved or nonconformal comparison with different structure-preservation requirements.

\begin{example}[The nonsimply laced central coefficient]
\label{ex:cbc26-b2}
For $B_2$, take an orthonormal basis $e_1,e_2$ and positive roots $e_1-e_2,e_2,e_1,e_1+e_2$.
The coroot of $\alpha$ is $2\alpha/(\alpha,\alpha)$.
Thus
\[
 \rho=(3/2,1/2),\qquad\rho^\vee=(2,1),\qquad
 (\rho,\rho^\vee)=7/2.
\]
The principal charge is $c(t)=86-60t-30/t$, so $c(t)+c(-t)=172$.
The rank is $2$, the dimension is $2+2\cdot4=10$, and $h^\vee=1+(\rho,e_1+e_2)=3$.
Consequently $2r+4h^\vee\dim\mathfrak g=124$ cannot be the reflected sum for all simple Lie algebras.
For simply laced root systems $\rho=\rho^\vee$; the general expression remains $2r+48(\rho,\rho^\vee)$.
\end{example}

\begin{proposition}[Two genus-one scalar normalizations]
\label{prop:cbc26-genus-one-normalization}
Let $B_2=1/6$ be the second Bernoulli number and define the scalar coefficient
\[
 F_1(\kappa)=\kappa\frac{2^1-1}{2^1}\frac{|B_2|}{2!}.
\]
Then $F_1(\kappa)=\kappa/24$.
Under $\kappa=C/2$, this is $C/48$.
The vacuum-energy exponent in $q^{L_0-C/24}$ is instead $-C/24$.
At $C=26$ these values are $13/24$ and $-13/12$.
\end{proposition}
\begin{proof}
The Bernoulli coefficient is $(1/2)(1/6)(1/2)=1/24$.
Substitution gives the other values.
\end{proof}

This proposition concerns the specified Bernoulli coefficient.
Its identification with a Hodge integral or a trace requires that map and its normalization.
The vacuum energy uses the separately defined conformal trace.

\begin{theorem}[A polynomial Batalin--Vilkovisky differential and its quotient]
\label{thm:cbc26-finite-bv}
Let $R=\mathbb C[x_1,\ldots,x_n]$ and $S\in R$.
Give $K=R\otimes\Lambda(\eta_1,\ldots,\eta_n)$ degrees $|x_i|=0$, $|\eta_i|=-1$.
Use the degree-one odd Poisson bracket with $\{x_i,\eta_j\}=\delta_{ij}$, $\{x_i,x_j\}=\{\eta_i,\eta_j\}=0$, extended by graded skew symmetry and the graded Leibniz rules.
Then $Q=\{S,-\}$ is the Koszul differential
\[
 Qx_i=0,\qquad Q\eta_i=\partial_i S,
\]
and satisfies $Q^2=0$.
If $\partial_1S,\ldots,\partial_nS$ is a regular sequence, the chain algebra projection
\[
 p:K\longrightarrow R/(\partial_1S,\ldots,\partial_nS),
 \qquad p(\eta_i)=0
\]
is a quasi-isomorphism.
\end{theorem}
\begin{proof}
The bracket formulas give the displayed derivation.
Its square vanishes on every generator and hence everywhere.
The projection is a chain map because it kills each $Q\eta_i$.
For a single regular element $f$, the complex $R\eta\xrightarrow{f}R$ has zero cohomology in degree $-1$ and quotient $R/(f)$ in degree zero.
Inductively tensor with the next two-term complex.
The two-column spectral sequence first computes the preceding Koszul cohomology, then multiplication by the next element on the quotient.
Regularity makes that multiplication injective.
The filtration is finite, so induction proves concentration in degree zero and the asserted projection isomorphism.
\end{proof}

For $S=\tfrac12\sum_i x_i^2$, the target is $\mathbb C$ and the projection evaluates $x_i=\eta_i=0$.
This is an explicit finite polynomial BRST comparison.
It has no collision sheaves, ultraviolet completion, or curvewise descent.
A comparison with chiral chains must construct those additional objects and maps.

\section{Form degree in higher interaction terms}

\begin{proposition}[Quartic expressions on a three-manifold]
\label{prop:cbc26-quartic-degree}
Let $M$ be a smooth three-manifold and $A\in\Omega^1(M;W)$ for a finite-dimensional vector space $W$.
Any expression obtained by applying a four-linear map to coefficients of $A$ and wedging all four one-form factors is zero.
If a vector field $v$ and a four-linear map $T:W^{\otimes4}\to\mathbb C$ are specified, the expression
\[
 T(\iota_vA,A,A,A)\in\Omega^3(M)
\]
is instead a defined three-form.
It can be nonzero, but depends on the additional contraction and tensor.
\end{proposition}
\begin{proof}
The first expression belongs to $\Omega^4(M)=0$.
For the second, $\iota_v A$ has form degree zero, so the total degree is three.
For an explicit local example take $W=\mathbb C^3$, $A=e_1dx+e_2dy+e_3dz$, $v=\partial_x$, and the tensor whose only nonzero coordinate entry is $T(e_1,e_1,e_2,e_3)=1$.
Then the expression is $dx\wedge dy\wedge dz$.
\end{proof}

An RTT relation is an algebraic relation on matrix-valued series.
To use it as an action term, a map to a three-form such as the displayed contraction must be defined and its required symmetries proved.
A critical affine level alone does not make the Lie bracket in $\Omega^\bullet(M)\otimes\mathfrak g$ vanish.
For example, $A=e\,dx+f\,dy$ in $\mathfrak{sl}_2$ has $[A,A]/2=h\,dx\wedge dy\ne0$, independently of any level parameter.

\begin{proposition}[Integration on a finite quotient]
\label{prop:cbc26-quotient-integral}
Let a finite group $\Gamma$ act freely by orientation-preserving diffeomorphisms on a closed oriented three-manifold $M$.
Let $\omega$ be an invariant three-form and $\bar\omega$ its descended form on $M/\Gamma$.
Then
\[
 \int_{M/\Gamma}\bar\omega=\frac1{|\Gamma|}\int_M\omega.
\]
\end{proposition}
\begin{proof}
The quotient map is an oriented covering of degree $|\Gamma|$.
An evenly covered coordinate cover and a partition of unity reduce integration to the sum over its $|\Gamma|$ sheets.
Each sheet has the same integral, giving the formula.
\end{proof}

For a nonfree finite action, integration on the quotient stack uses the right-hand side as its normalization.
An additional division by $|\Gamma|$ after taking this quotient integral changes the normalization.
Neither a fractional grading nor a W-algebra specifies an action of $\Gamma$ on $M$.
